\documentclass[11pt,a4paper]{article}

% Packages
\usepackage[utf8]{inputenc}
\usepackage[T1]{fontenc}
\usepackage{amsmath,amssymb,amsfonts}
\usepackage[margin=1in]{geometry}
\usepackage{enumitem}
\usepackage{hyperref}
\usepackage{fancyhdr}
\usepackage{titlesec}
\usepackage{tocloft}

% Page style
\pagestyle{fancy}
\fancyhf{}
\fancyhead[L]{\small Lecture Notes: Induction and Recursion}
\fancyhead[R]{\small \thepage}
\renewcommand{\headrulewidth}{0.4pt}

% Section formatting
\titleformat{\section}{\Large\bfseries}{Section \thesection}{1em}{}
\titleformat{\subsection}{\large\bfseries}{\thesubsection}{1em}{}

% Adjust table of contents appearance
\renewcommand{\cftsecfont}{\bfseries}
\renewcommand{\cftsubsecfont}{\normalfont}

\begin{document}

% Title page
\begin{titlepage}
    \centering
    \vspace*{4cm}
    
    {\Huge\bfseries Lecture Notes:\\ Induction and Recursion\par}
    \vspace{1cm}
    {\Large Discrete Mathematics\par}
    \vspace{2cm}
    {\large \textbf{Author:} Bakdaulet Sotsial\par}
    \vspace{1cm}
    
    \vfill
    {\small Astana IT University \par}
    \vspace*{2cm}
\end{titlepage}

% Table of contents
\tableofcontents
\newpage

% Section 1: Mathematical Induction
\section{Mathematical Induction}

\subsection{Principle of Mathematical Induction}

\textbf{Theorem (Principle of Mathematical Induction).} Let $P(n)$ be a propositional function defined for integers $n \geq n_0$, where $n_0$ is a fixed integer (usually $n_0 = 0$ or $1$). Suppose that:
\begin{enumerate}[label=(\roman*)]
    \item \textbf{Base case:} $P(n_0)$ is true.
    \item \textbf{Inductive step:} For every integer $k \geq n_0$, if $P(k)$ is true, then $P(k+1)$ is true.
\end{enumerate}
Then $P(n)$ is true for all integers $n \geq n_0$.

\textbf{Remark.} Mathematical induction is a fundamental proof technique for statements involving natural numbers. It formalizes the "domino effect": if the first domino falls, and each domino knocks over the next one, then all dominoes will eventually fall.

\subsection{Structure of an Inductive Proof}

A proof by mathematical induction consists of the following essential components:

\begin{enumerate}[label=\textbf{\arabic*.}]
    \item \textbf{Define the predicate:} Clearly state $P(n)$, the statement to be proved for each $n \geq n_0$.
    
    \item \textbf{Base case:} Verify that $P(n_0)$ holds. This anchors the induction.
    
    \item \textbf{Inductive hypothesis:} Assume that $P(k)$ is true for some arbitrary $k \geq n_0$.
    
    \item \textbf{Inductive step:} Using the inductive hypothesis, prove that $P(k+1)$ is true.
    
    \item \textbf{Conclusion:} By the principle of mathematical induction, $P(n)$ holds for all $n \geq n_0$.
\end{enumerate}

\subsection{Examples of Proofs by Induction}

\begin{example}
Prove that for all $n \geq 1$,
\begin{equation*}
    1 + 2 + 3 + \cdots + n = \frac{n(n+1)}{2}.
\end{equation*}

\textbf{Proof.} Let $P(n)$ be the statement: $\displaystyle\sum_{i=1}^{n} i = \frac{n(n+1)}{2}$.

\textit{Base case ($n = 1$):} Left side: $1$. Right side: $\frac{1 \cdot 2}{2} = 1$. So $P(1)$ is true.

\textit{Inductive hypothesis:} Assume $P(k)$ is true for some $k \geq 1$, i.e.,
\begin{equation*}
    1 + 2 + \cdots + k = \frac{k(k+1)}{2}.
\end{equation*}

\textit{Inductive step:} We must prove $P(k+1)$:
\begin{equation*}
    1 + 2 + \cdots + k + (k+1) = \frac{(k+1)(k+2)}{2}.
\end{equation*}
Starting from the left side and using the inductive hypothesis:
\begin{align*}
    1 + 2 + \cdots + k + (k+1) &= \frac{k(k+1)}{2} + (k+1) \\
    &= (k+1)\left(\frac{k}{2} + 1\right) \\
    &= (k+1)\left(\frac{k+2}{2}\right) \\
    &= \frac{(k+1)(k+2)}{2}.
\end{align*}
Thus $P(k+1)$ holds.

\textit{Conclusion:} By mathematical induction, $P(n)$ is true for all $n \geq 1$.
\end{example}

\begin{example}
Prove that $n! > 2^n$ for all integers $n \geq 4$.

\textbf{Proof.} Let $P(n)$ be: $n! > 2^n$.

\textit{Base case ($n = 4$):} $4! = 24$, $2^4 = 16$, and $24 > 16$. So $P(4)$ holds.

\textit{Inductive hypothesis:} Assume $P(k)$ for some $k \geq 4$, i.e., $k! > 2^k$.

\textit{Inductive step:} Show $(k+1)! > 2^{k+1}$.
\begin{align*}
    (k+1)! &= (k+1) \cdot k! \\
    &> (k+1) \cdot 2^k \quad \text{(by inductive hypothesis)}\\
    &\geq 5 \cdot 2^k \quad \text{(since } k \geq 4, \, k+1 \geq 5\text{)} \\
    &> 2 \cdot 2^k = 2^{k+1}.
\end{align*}
Thus $P(k+1)$ holds. By induction, $n! > 2^n$ for all $n \geq 4$.
\end{example}

\subsection{Strong Induction}

\textbf{Theorem (Principle of Strong Induction).} Let $P(n)$ be a propositional function for $n \geq n_0$. Suppose that:
\begin{enumerate}[label=(\roman*)]
    \item \textbf{Base case:} $P(n_0)$ is true.
    \item \textbf{Strong inductive step:} For every $k \geq n_0$, if $P(n_0), P(n_0+1), \dots, P(k)$ are all true, then $P(k+1)$ is true.
\end{enumerate}
Then $P(n)$ is true for all $n \geq n_0$.

\textbf{Remark.} Strong induction differs from ordinary induction in the hypothesis: we assume the truth of $P$ for \textit{all} previous values, not just the immediate predecessor. This is especially useful when $P(k+1)$ depends on multiple earlier values, not just $P(k)$.

\begin{example}
Prove that every integer $n \geq 2$ is either prime or can be written as a product of primes.

\textbf{Proof (by strong induction).}
\textit{Base case ($n = 2$):} $2$ is prime, so the statement holds.

\textit{Inductive hypothesis:} Assume the statement holds for all integers $m$ with $2 \leq m \leq k$ for some $k \geq 2$.

\textit{Inductive step:} Consider $k+1$. If $k+1$ is prime, we are done. If not, then $k+1 = ab$ where $2 \leq a, b \leq k$. By the inductive hypothesis, both $a$ and $b$ are either prime or products of primes. Hence $k+1$ is a product of primes. This completes the proof.
\end{example}

% Section 2: Recursively Defined Sequences
\section{Recursively Defined Sequences}

\textbf{Definition (Recursion).} A sequence $\{a_n\}$ is said to be \textit{recursively defined} (or defined by \textit{recurrence}) if each term is defined in terms of previous terms. A recursive definition consists of two parts:
\begin{enumerate}[label=(\roman*)]
    \item \textbf{Initial condition(s):} The value(s) of the first one or more terms, given explicitly.
    \item \textbf{Recurrence relation:} An equation that expresses $a_n$ in terms of earlier terms $a_{n-1}, a_{n-2}, \dots$ for all $n$ beyond the initial conditions.
\end{enumerate}

\textbf{Remark.} Recursive definitions are ubiquitous in mathematics and computer science. They allow complex objects to be built from simpler ones, and they naturally lead to inductive proofs of properties.

\begin{example}[Factorial]
The factorial function $n!$ can be defined recursively:
\begin{align*}
    0! &= 1 \quad \text{(initial condition)}, \\
    n! &= n \cdot (n-1)! \quad \text{for } n \geq 1 \quad \text{(recurrence relation)}.
\end{align*}
\end{example}

\begin{example}[Arithmetic Progression]
An arithmetic progression with first term $a$ and common difference $d$:
\begin{align*}
    a_0 &= a, \\
    a_n &= a_{n-1} + d \quad \text{for } n \geq 1.
\end{align*}
\end{example}

\begin{example}[Geometric Progression]
A geometric progression with first term $a$ and common ratio $r$:
\begin{align*}
    a_0 &= a, \\
    a_n &= r \cdot a_{n-1} \quad \text{for } n \geq 1.
\end{align*}
\end{example}

\begin{example}[Fibonacci Sequence]
The famous Fibonacci sequence $\{F_n\}_{n=0}^{\infty}$ is defined recursively:
\begin{align*}
    F_0 &= 0, \quad F_1 = 1 \quad \text{(initial conditions)}, \\
    F_n &= F_{n-1} + F_{n-2} \quad \text{for } n \geq 2 \quad \text{(recurrence relation)}.
\end{align*}
The first few terms are: $0, 1, 1, 2, 3, 5, 8, 13, 21, 34, 55, \dots$

\textbf{Interpretation:} Each Fibonacci number (after the first two) is the sum of the two preceding numbers. The sequence appears in numerous natural phenomena, algorithm analysis, and combinatorial problems.
\end{example}

% Section 3: Recurrence Relations
\section{Recurrence Relations}

\textbf{Definition (Recurrence Relation).} A \textit{recurrence relation} for a sequence $\{a_n\}$ is an equation that expresses $a_n$ as a function of one or more of the preceding terms $a_{n-1}, a_{n-2}, \dots, a_{n-k}$ for all $n \geq k$, where $k$ is a fixed integer. The relation must hold for all $n$ beyond some initial index.

\textbf{Definition (Order of a Recurrence).} The \textit{order} of a recurrence relation is the number of preceding terms on which $a_n$ depends. For example, $a_n = a_{n-1} + a_{n-2}$ has order $2$.

\subsection{Types of Recurrence Relations}

\begin{enumerate}[label=(\arabic*)]
    \item \textbf{Linear vs. Nonlinear:} A recurrence is \textit{linear} if it is a linear combination of previous terms. Otherwise, it is nonlinear.
    
    \textbf{Example:} $a_n = 3a_{n-1} - 2a_{n-2}$ is linear. $a_n = a_{n-1}^2 + a_{n-2}$ is nonlinear.
    
    \item \textbf{Homogeneous vs. Non-homogeneous:} A linear recurrence is \textit{homogeneous} if the linear combination of previous terms equals zero when we bring all terms involving $a$ to one side. It is \textit{non-homogeneous} if there is an additional term that is a function of $n$ alone.
    
    \textbf{Example:} $a_n = 5a_{n-1} - 6a_{n-2}$ is homogeneous.
    $a_n = 5a_{n-1} - 6a_{n-2} + 2^n$ is non-homogeneous.
    
    \item \textbf{Constant vs. Variable Coefficients:} If the coefficients multiplying the previous terms are constants (independent of $n$), we say the recurrence has \textit{constant coefficients}. If they depend on $n$, they are \textit{variable coefficients}.
    
    \textbf{Example:} $a_n = 3a_{n-1} - 2a_{n-2}$ has constant coefficients.
    $a_n = n a_{n-1} + a_{n-2}$ has variable coefficients.
\end{enumerate}

\textbf{Definition (Linear Homogeneous Recurrence with Constant Coefficients).} A \textit{linear homogeneous recurrence relation with constant coefficients} of order $k$ has the form:
\begin{equation*}
    a_n = c_1 a_{n-1} + c_2 a_{n-2} + \cdots + c_k a_{n-k},
\end{equation*}
where $c_1, c_2, \dots, c_k \in \mathbb{R}$ are constants and $c_k \neq 0$. This is the main class of recurrence relations for which systematic solution methods exist.

\begin{example}
The Fibonacci recurrence $F_n = F_{n-1} + F_{n-2}$ is a linear homogeneous recurrence with constant coefficients of order $2$ (here $c_1 = 1$, $c_2 = 1$).
\end{example}

\begin{example}
The recurrence $a_n = 2a_{n-1} - a_{n-2}$ is a linear homogeneous recurrence with constant coefficients of order $2$.
\end{example}

\begin{example}[Non-homogeneous]
$a_n = a_{n-1} + a_{n-2} + n$ is linear but non-homogeneous due to the $+n$ term.
\end{example}

% Section 4: Solving Recurrence Relations
\section{Solving Recurrence Relations}

We focus on solving \textit{linear homogeneous recurrence relations with constant coefficients}. The goal is to find an explicit formula (a \textit{closed form}) for $a_n$ that depends only on $n$ and the initial conditions.

\subsection{Method of Characteristic Polynomial}

Consider a linear homogeneous recurrence relation of order $k$:
\begin{equation*}
    a_n = c_1 a_{n-1} + c_2 a_{n-2} + \cdots + c_k a_{n-k}.
\end{equation*}
We look for solutions of the form $a_n = r^n$, where $r \neq 0$ is a constant to be determined.

Substituting $a_n = r^n$ into the recurrence:
\begin{align*}
    r^n &= c_1 r^{n-1} + c_2 r^{n-2} + \cdots + c_k r^{n-k} \\
    r^n - c_1 r^{n-1} - c_2 r^{n-2} - \cdots - c_k r^{n-k} &= 0 \\
    r^{n-k}\bigl(r^k - c_1 r^{k-1} - c_2 r^{k-2} - \cdots - c_k\bigr) &= 0.
\end{align*}
Since $r \neq 0$, we obtain the \textbf{characteristic equation}:
\begin{equation*}
    r^k - c_1 r^{k-1} - c_2 r^{k-2} - \cdots - c_k = 0.
\end{equation*}
The roots of this polynomial are called the \textit{characteristic roots}.

\subsection{Case 1: Distinct Real Roots}

\textbf{Theorem.} If the characteristic equation has $k$ distinct real roots $r_1, r_2, \dots, r_k$, then the general solution to the recurrence relation is:
\begin{equation*}
    a_n = \alpha_1 r_1^n + \alpha_2 r_2^n + \cdots + \alpha_k r_k^n,
\end{equation*}
where $\alpha_1, \alpha_2, \dots, \alpha_k$ are constants determined by the initial conditions.

\begin{example}
Solve the recurrence $a_n = 5a_{n-1} - 6a_{n-2}$ with initial conditions $a_0 = 2$, $a_1 = 5$.

\textbf{Solution.}
The recurrence can be rewritten as $a_n - 5a_{n-1} + 6a_{n-2} = 0$.
Characteristic equation: $r^2 - 5r + 6 = 0$.
Factor: $(r-2)(r-3) = 0$, so $r_1 = 2$, $r_2 = 3$ (distinct).
General solution: $a_n = \alpha_1 \cdot 2^n + \alpha_2 \cdot 3^n$.

Use initial conditions:
\begin{align*}
    n = 0 &: \alpha_1 + \alpha_2 = 2, \\
    n = 1 &: 2\alpha_1 + 3\alpha_2 = 5.
\end{align*}
Solving: subtract twice the first equation from the second: $(2\alpha_1 + 3\alpha_2) - 2(\alpha_1 + \alpha_2) = 5 - 4 \implies \alpha_2 = 1$. Then $\alpha_1 = 2 - 1 = 1$.
Thus, $a_n = 2^n + 3^n$.
\end{example}

\begin{example}[Fibonacci Numbers]
Solve the Fibonacci recurrence $F_n = F_{n-1} + F_{n-2}$ with $F_0 = 0$, $F_1 = 1$.

\textbf{Solution.} Rewrite: $F_n - F_{n-1} - F_{n-2} = 0$.
Characteristic equation: $r^2 - r - 1 = 0$.
Roots:
\begin{equation*}
    r_{1,2} = \frac{1 \pm \sqrt{1 + 4}}{2} = \frac{1 \pm \sqrt{5}}{2}.
\end{equation*}
Let $\varphi = \frac{1+\sqrt{5}}{2}$ (the golden ratio) and $\psi = \frac{1-\sqrt{5}}{2}$.
General solution: $F_n = \alpha_1 \varphi^n + \alpha_2 \psi^n$.

Using initial conditions:
\begin{align*}
    n = 0 &: \alpha_1 + \alpha_2 = 0, \\
    n = 1 &: \alpha_1 \varphi + \alpha_2 \psi = 1.
\end{align*}
From the first, $\alpha_2 = -\alpha_1$. Substitute into the second:
\begin{equation*}
    \alpha_1 \varphi - \alpha_1 \psi = 1 \implies \alpha_1 (\varphi - \psi) = 1.
\end{equation*}
Now, $\varphi - \psi = \frac{1+\sqrt{5}}{2} - \frac{1-\sqrt{5}}{2} = \sqrt{5}$.
So $\alpha_1 = \frac{1}{\sqrt{5}}$, $\alpha_2 = -\frac{1}{\sqrt{5}}$.

Thus, \textbf{Binet's formula}:
\begin{equation*}
    F_n = \frac{1}{\sqrt{5}}\left(\frac{1+\sqrt{5}}{2}\right)^n - \frac{1}{\sqrt{5}}\left(\frac{1-\sqrt{5}}{2}\right)^n.
\end{equation*}
\end{example}

\subsection{Case 2: Repeated Real Roots}

\textbf{Theorem.} If the characteristic equation has a root $r$ of multiplicity $m$ (i.e., $(x - r)^m$ is a factor), then the terms $r^n$, $n r^n$, $n^2 r^n$, $\dots$, $n^{m-1} r^n$ are all solutions. The part of the general solution corresponding to this root is:
\begin{equation*}
    (\alpha_1 + \alpha_2 n + \alpha_3 n^2 + \cdots + \alpha_m n^{m-1}) r^n.
\end{equation*}
If there are several repeated roots, the general solution is the sum of such expressions for each distinct root.

\begin{example}
Solve $a_n = 6a_{n-1} - 9a_{n-2}$ with $a_0 = 1$, $a_1 = 6$.

\textbf{Solution.} Rewrite: $a_n - 6a_{n-1} + 9a_{n-2} = 0$.
Characteristic equation: $r^2 - 6r + 9 = 0$, or $(r-3)^2 = 0$.
Double root $r = 3$ (multiplicity 2).
General solution: $a_n = \alpha_1 \cdot 3^n + \alpha_2 n \cdot 3^n$.

Using initial conditions:
\begin{align*}
    n = 0 &: \alpha_1 = 1, \\
    n = 1 &: 3\alpha_1 + 3\alpha_2 = 6 \implies 3(1) + 3\alpha_2 = 6 \implies 3\alpha_2 = 3 \implies \alpha_2 = 1.
\end{align*}
Thus, $a_n = 3^n + n \cdot 3^n = (n+1) \cdot 3^n$.
\end{example}

\begin{example}[Triple Root]
Consider $a_n = 3a_{n-1} - 3a_{n-2} + a_{n-3}$.
Characteristic equation: $r^3 - 3r^2 + 3r - 1 = 0$, or $(r-1)^3 = 0$.
Root $r = 1$ with multiplicity 3.
General solution: $a_n = \alpha_1 \cdot 1^n + \alpha_2 n \cdot 1^n + \alpha_3 n^2 \cdot 1^n = \alpha_1 + \alpha_2 n + \alpha_3 n^2$.
\end{example}

\subsection{Summary of the Method}

\textbf{Algorithm for Solving Linear Homogeneous Recurrence Relations with Constant Coefficients:}

\begin{enumerate}[label=\textbf{Step \arabic*:}]
    \item Write the recurrence in standard form:
    \begin{equation*}
        a_n - c_1 a_{n-1} - c_2 a_{n-2} - \cdots - c_k a_{n-k} = 0.
    \end{equation*}
    \item Form the \textbf{characteristic equation}:
    \begin{equation*}
        r^k - c_1 r^{k-1} - c_2 r^{k-2} - \cdots - c_k = 0.
    \end{equation*}
    \item Find all roots $r_1, r_2, \dots$ of the characteristic equation along with their multiplicities.
    \item Write the general solution:
    \begin{itemize}
        \item For each distinct root $r$ of multiplicity $m$, include terms:
        \begin{equation*}
            (\alpha_1 + \alpha_2 n + \cdots + \alpha_m n^{m-1}) r^n.
        \end{equation*}
    \end{itemize}
    \item Use the initial conditions to set up a system of linear equations and solve for the constants $\alpha_i$.
\end{enumerate}

\textbf{Remark.} The method described here applies specifically to linear homogeneous recurrences with constant coefficients. Non-homogeneous recurrences require finding a particular solution (similar to the method of undetermined coefficients for differential equations), and recurrences with variable coefficients often require different techniques (e.g., generating functions).

\textbf{Applications.} Recurrence relations and their solutions are fundamental in analyzing the time complexity of recursive algorithms (e.g., divide-and-conquer algorithms like MergeSort), modeling population growth, computing combinatorial quantities (e.g., Catalan numbers), and many other areas in mathematics and computer science.

\end{document}